% Edition check: concise; total-page review limit 10, including references.
% Reader task: Answer one coherent thematic question in an independently readable account, not a concatenation of Parts I–IV.
% Select one complete thread; remove genuine side branches before drafting.
% In the Introduction state every selected principal conclusion, not only names.
% Review body-to-introduction reverse coverage after changing any result or scope.
% Style: delete trivial transitions; retain useful self-contained repetition;
% use itemize/enumerate; put secondary consequences afterward.
% An introduction consisting of definitions and theorems is valid:
% definitions and theorems are valid without a narrative roadmap.
% Two slots illustrate wiring, not a theorem-count requirement. Remove unused slots.
\documentclass[12pt,reqno]{amsart}
\usepackage{math-review}
\title{\placeholder{Topic}: A Thematic Review}
\author{}
\date{\placeholder{Checked literature cutoff}}
\newcommand{\PrincipalStatement}{%
\placeholder{Specify objects, all necessary hypotheses, quantifiers and ranges.}
\begin{itemize}
\item \placeholder{A necessary condition; omit this list if there is no real list.}
\item \placeholder{A second independent condition, only when needed.}
\end{itemize}
\placeholder{State the actual mathematical conclusion and exact source/status.}
}
\newcommand{\SecondStatement}{%
\placeholder{State another selected principal answer with its own data, necessary conditions and actual conclusion; remove this slot when outside scope.}
}
% Optional unresolved-core slot, not a Problem quota. Remove it when the scope
% has no selected unresolved question; explain any settled organizing question instead.
\newcommand{\CoreProblemStatement}{%
\placeholder{State the core unresolved target with its objects, assumptions, quantifiers and exact conclusion to prove. Use the sourced formulation, not a vague research direction.}
}
\begin{document}
\begin{abstract}
\placeholder{Give the question and selected answers or boundary in two or three sentences. The abstract does not replace the main statements below.}
\end{abstract}
\maketitle

\section{Introduction}\label{sec:introduction}
\placeholder{Define the minimum necessary notation. State fixed data, unknowns and branch conditions before using them. Give actual answers with quantifiers, parameter ranges and source/status qualifications. Delete transitions that merely announce what follows.}
\begin{theorem}\label{thm:principal}
\PrincipalStatement
\end{theorem}
\begin{theorem}\label{thm:second}
\SecondStatement
\end{theorem}
\begin{problem}\label{prob:core-target}
\CoreProblemStatement
\end{problem}
\placeholder{Keep a precise remaining question or negative boundary here when it is a principal commitment. Change theorem to proposition, sourced claim or Problem as mathematically appropriate. No invented optional titles.}

\section{Results through their mechanisms}\label{sec:mechanism}
\begin{theoremrecall}{thm:principal}
\PrincipalStatement
\end{theoremrecall}
\placeholder{Develop the selected results through their decisive mechanisms and boundaries. Retain definitions at first use. Do not import every proof from III or replace the article by a theorem catalog. A recall is optional if it adds backtracking rather than removing it. Reuse the canonical body whenever a full repetition is retained.}

\section{The second answer and the boundary}\label{sec:boundary}
\begin{theoremrecall}{thm:second}
\SecondStatement
\end{theoremrecall}
\placeholder{Use a mathematical heading. Explain the second answer or boundary without bundling secondary consequences into its main statement. Define a precise sourced Problem only if the selected scope needs one. Do not manufacture a fixed number of Problems.}

\placeholder{For Problem~\ref{prob:core-target}, give its primary source, checked cutoff status, strongest relevant solved range and exact remaining range. Do not infer openness from a failed search. Remove this paragraph with the optional Problem slot when it is outside the selected scope.}

\templatereferences
\end{document}
